% GENERATED FILE. Edit canonical lessons, then run npm run generate:books.
% canonical-source-hash: fnv1a64:78e9caf18c2ac3ca
% canonical-lessons: TE-C71-uncle, TE-C71-aunt, TE-C71-brother-in-law, TE-C71-sister-in-law, TE-C71-son-in-law

\chapter{The Wider Family}
\label{ch:te-the-wider-family}
\hlchaptermodality{71}

\begin{chapteropening}
\textbf{By the end of this chapter:} \emph{I can name the family a Telugu marriage adds, and ask after each of them.}

Complete TE-C71-son-in-law and demonstrate the chapter promise.
\end{chapteropening}

\section[māma]{\te{మామ} (māma) — an uncle; a father-in-law}

\label{lesson:TE-C71-uncle}



\begin{quote}
Before the new one: what did \emph{sālīḍu} mean?
\end{quote}

\subsection*{You'll want to know: \te{మామ}}
\textbf{\te{మామ}} (\emph{māma}) — \textquotedblleft{}an uncle; a father-in-law\textquotedblright{}.

\te{మామ} is a mother's brother, and it is also a father-in-law. Telugu keeps one word for both because for a very long time they were often the same man.

That is the whole of it: a man's sister's daughter was the expected match for his son, so the uncle a child grew up with became the father-in-law they married into.

Asking after a \te{మామ}'s \te{ఆరోగ్యం} is the first thing said after \te{నమస్కారం} in a house you have married into.

\begin{grammarlens}[title={What you've built}]
One. And the reason Telugu needs fewer family words than you expect.
\end{grammarlens}

\subsection*{Your turn}
\begin{itemize}
\raggedright
  \item \emph{Say it:} \emph{māma}
  \item \emph{Say it:} it once more, slowly
  \item \emph{Say it:} \emph{māma}, then \emph{ārōgyaṁ}, as one question
  \item \emph{Recall:} say \emph{goḍḍali}, then read \textbf{\te{ఎలుక}} and say what it means
\end{itemize}

\begin{tcolorbox}[breakable,colback=teal!4,colframe=teal!35!black,title={Before you move on}]
What does \te{మామ} mean? (\textquotedblleft{}an uncle; a father-in-law\textquotedblright{}.) Say it once more.
\end{tcolorbox}
\section[atta]{\te{అత్త} (atta) — an aunt; a mother-in-law}

\label{lesson:TE-C71-aunt}



\begin{quote}
Before the new one: what did \emph{māma} mean?
\end{quote}

\subsection*{You'll want to know: \te{అత్త}}
\textbf{\te{అత్త}} (\emph{atta}) — \textquotedblleft{}an aunt; a mother-in-law\textquotedblright{}.

\te{అత్త} is the \te{మామ}'s counterpart: a father's sister, and a mother-in-law.

The same reasoning folds them together. She is the woman a child calls \te{అత్త} while growing up and the woman they answer to after marrying, and one word covers both lives.

It is inherited and it is among the oldest words in the family. The neighbouring languages all keep something very close to it.

\begin{grammarlens}[title={What you've built}]
Two: \te{మామ} and \te{అత్త}, and they usually arrive as a pair.
\end{grammarlens}

\subsection*{Your turn}
\begin{itemize}
\raggedright
  \item \emph{Say it:} \emph{atta}
  \item \emph{Say it:} it once more, slowly
  \item \emph{Say it:} \emph{māma}, then \emph{atta}, and say which one is the mother-in-law
  \item \emph{Recall:} read \textbf{\te{రంపం}}, then say \emph{kappa}
\end{itemize}

\begin{tcolorbox}[breakable,colback=teal!4,colframe=teal!35!black,title={Before you move on}]
What does \te{అత్త} mean? (\textquotedblleft{}an aunt; a mother-in-law\textquotedblright{}.) Say it once more.
\end{tcolorbox}
\section[bāva]{\te{బావ} (bāva) — a brother-in-law; a male cousin}

\label{lesson:TE-C71-brother-in-law}



\begin{quote}
Before the new one: what did \emph{atta} mean?
\end{quote}

\subsection*{You'll want to know: \te{బావ}}
\textbf{\te{బావ}} (\emph{bāva}) — \textquotedblleft{}a brother-in-law; a male cousin\textquotedblright{}.

\te{బావ} is a sister's husband, a wife's brother, and the \te{మామ}'s son — the cousin you might have married.

So the word covers a bond that is close, level and a little teasing. Two men who call each other \te{బావ} can joke in a way an \te{అన్న} and a \te{తమ్ముడు} cannot.

City Telugu now uses it for a friend, roughly where English says mate. That irritates some people and has stopped nobody.

\begin{grammarlens}[title={What you've built}]
Three. The first one standing on level ground with you.
\end{grammarlens}

\subsection*{Your turn}
\begin{itemize}
\raggedright
  \item \emph{Say it:} \emph{bāva}
  \item \emph{Say it:} it once more, slowly
  \item \emph{Say it:} \emph{bāva}, then \emph{māma}, and say which one is older
  \item \emph{Recall:} say \emph{pāra}, then read \textbf{\te{దోమ}} and say what it means
\end{itemize}

\begin{tcolorbox}[breakable,colback=teal!4,colframe=teal!35!black,title={Before you move on}]
What does \te{బావ} mean? (\textquotedblleft{}a brother-in-law; a male cousin\textquotedblright{}.) Say it once more.
\end{tcolorbox}
\section[maradalu]{\te{మరదలు} (maradalu) — a sister-in-law; a female cousin}

\label{lesson:TE-C71-sister-in-law}



\begin{quote}
Before the new one: what did \emph{bāva} mean?
\end{quote}

\subsection*{You'll want to know: \te{మరదలు}}
\textbf{\te{మరదలు}} (\emph{maradalu}) — \textquotedblleft{}a sister-in-law; a female cousin\textquotedblright{}.

\te{మరదలు} is the \te{బావ}'s counterpart on the other side: a wife's younger sister, a younger brother's wife, and the \te{మామ}'s daughter — the cousin you might have married.

So it carries the same easy, teasing register that \te{బావ} does. A \te{బావ} and a \te{మరదలు} can joke across a room in a way an \te{అన్న} and a \te{చెల్లి} cannot, and Telugu films have lived off that for eighty years.

It is inherited, and it belongs to the same reckoning as the rest of this chapter: the person a marriage hands you was very often somebody the family already had a word for.

\begin{grammarlens}[title={What you've built}]
Four. The one who matches \te{బావ}.
\end{grammarlens}

\subsection*{Your turn}
\begin{itemize}
\raggedright
  \item \emph{Say it:} \emph{maradalu}
  \item \emph{Say it:} it once more, slowly
  \item \emph{Say it:} \emph{bāva}, then \emph{maradalu}, and say which of the two is the woman
  \item \emph{Recall:} read \textbf{\te{సుత్తి}}, then say \emph{uḍuta}
\end{itemize}

\begin{tcolorbox}[breakable,colback=teal!4,colframe=teal!35!black,title={Before you move on}]
What does \te{మరదలు} mean? (\textquotedblleft{}a sister-in-law; a female cousin\textquotedblright{}.) Say it once more.
\end{tcolorbox}
\section[alluḍu]{\te{అల్లుడు} (alluḍu) — a son-in-law}

\label{lesson:TE-C71-son-in-law}



\begin{quote}
Before the new one: what did \emph{maradalu} mean?
\end{quote}

\subsection*{You'll want to know: \te{అల్లుడు}}
\textbf{\te{అల్లుడు}} (\emph{alluḍu}) — \textquotedblleft{}a son-in-law\textquotedblright{}.

\te{అల్లుడు} is a daughter's husband — the \te{మామ}'s word for the young man on the far side of the match.

He is treated as a guest for life. An \te{అల్లుడు} on a visit is fed first, seated first and asked after first, and the household's own people wait their turn.

Which is what the chapter has been building to: \te{మామ}, \te{అత్త}, \te{బావ}, \te{మరదలు}, \te{అల్లుడు}. Ask each of them how they are and you have done the whole visit properly.

\begin{grammarlens}[title={What you've built}]
Five: \te{మామ}, \te{అత్త}, \te{బావ}, \te{మరదలు}, \te{అల్లుడు}. A marriage's worth of new people to ask after.
\end{grammarlens}

\subsection*{Your turn}
\begin{itemize}
\raggedright
  \item \emph{Say it:} \emph{alluḍu}
  \item \emph{Say it:} it once more, slowly
  \item \emph{Say it:} all five in order, then \emph{ārōgyaṁ}, asking after the last of them
  \item \emph{Recall:} say \emph{karra}, then read \textbf{\te{సాలీడు}} and say what it means
\end{itemize}

\begin{tcolorbox}[breakable,colback=teal!4,colframe=teal!35!black,title={Before you move on}]
What does \te{అల్లుడు} mean? (\textquotedblleft{}a son-in-law\textquotedblright{}.) Say it once more.
\end{tcolorbox}
